% Lösungen Einheit 1 von 3 – Modell und Glockenkurve: die Dichtefunktion lesen und skizzieren (μ als Symmetrieachse und Maximumsstelle, σ als Breitenmaß (zusammenbau v0.1)
\einheitenkopf{Einheit 1 von 3 · Modell und Glockenkurve: die Dichtefunktion lesen und skizzieren (μ als Symmetrieachse und Maximumsstelle, σ als Breitenmaß}

\erg{8}{a) Fläche unter der Dichte \quad b) Gleiche Symmetrieachse bei $50$; die Glocke wird doppelt so breit und halb so hoch – die Fläche bleibt eins. \quad c) $0$: Über einem einzelnen Wert gibt es keine Fläche; der Dichtewert an der Stelle $35$ ist keine Wahrscheinlichkeit. \quad d) Nein: $(500 - 0) : 4 = 125$ Standardabweichungen trennen $0$ vom Erwartungswert; $P(X < 0)$ ist verschwindend klein. Das Modell beschreibt die realistischen Füllmengen. \quad e) Vorschlag 1: dieselbe Glocke nach rechts verschoben (etwa $\mu = 1006$) – die Fläche links von $1000$ wird kleiner. Vorschlag 2: gleiche Mitte, schmaler und höher (die Fläche bleibt eins) – mehr Masse liegt nahe bei $\mu$, die Fläche links von $1000$ wird kleiner.}

8b)

\normalverteilung{50}{10}


8e)

\normalverteilung[bis=1000]{1006}{4} \normalverteilung[bis=1000]{1003}{2}

\erg{9}{a) Fehler: Der Dichtewert ist keine Wahrscheinlichkeit. Richtig: $P(X = 20) = 0$ – über einem Einzelwert gibt es keine Fläche. \quad b) Die Fläche über einem Bereich ist die Wahrscheinlichkeit, dass $X$ dort liegt. Über der ganzen Achse ist das das sichere Ereignis – Wahrscheinlichkeit eins.}
